\chapter{System Model}\label{ch:system}
PaperMind is organised as four layers over a single SQLite store (Figure~\ref{fig:arch}): an \emph{ingestion} layer that
turns documents into structured questions, an \emph{NLP analysis} layer that labels questions and computes historical
statistics, a \emph{GenAI} layer that plans, writes and verifies new papers, and an \emph{interaction} layer consisting
of a REST API, a web interface and a tool-using study assistant.

\tikzset{
  blk/.style={draw, rounded corners=3pt, align=center, font=\small, minimum height=0.9cm, fill=white, text width=#1},
  blk/.default=2.6cm,
  lay/.style={draw=gray!60, rounded corners=6pt, fill=gray!8, inner sep=6pt},
  arr/.style={-{Stealth[length=2.2mm]}, thick},
}

\begin{figure}[t]
\centering
\begin{tikzpicture}[x=3.45cm, y=1.55cm]
  \node[blk=2.5cm] (pdf)  at (0,0)  {Past papers\\(PDF / scan / photo)};
  \node[blk=2.5cm] (ing)  at (1,0)  {Ingestion\\text, OCR, parsing};
  \node[blk=2.5cm] (lab)  at (2,0)  {Labelling (LLM)\\topic, Bloom level};
  \node[blk=2.5cm] (db)   at (3,0)  {SQLite\\corpus};
  \node[blk=2.5cm] (ana)  at (2,-1.6) {Historical analysis\\weightage, hot topics};
  \node[blk=2.5cm] (gen)  at (3,-1.6) {Paper generator\\plan, RAG, verify};
  \node[blk=2.5cm] (asst) at (1,-1.6) {Study assistant\\(tool-using agent)};
  \node[blk=2.5cm] (ui)   at (0,-1.6) {Web UI, CLI,\\REST API};
  \node[blk=2.5cm] (stu)  at (0,-3.1) {Student};
  \node[blk=2.5cm, fill=blue!6] (llm) at (2,-3.1) {Google Gemini\\(LLM)};
  \draw[arr] (pdf) -- (ing); \draw[arr] (ing) -- (lab); \draw[arr] (lab) -- (db);
  \draw[arr] (db) -- (gen); \draw[arr] (db) -- (ana);
  \draw[arr] (ana) -- (asst);
  \draw[arr] (gen.south) -- ++(0,-0.35) -| ([xshift=6pt]asst.south);
  \draw[arr] (asst) -- (ui); \draw[arr, <->] (ui) -- (stu);
  \draw[arr, <->, dashed] (llm.east) -| (gen.south east);
  \draw[arr, <->, dashed] (llm.west) -| ([xshift=-6pt]asst.south);
\end{tikzpicture}
\caption{System model of PaperMind (dashed: calls to the LLM; labelling also uses it)}
\label{fig:arch}
\end{figure}

\section{Dataset Collection}
The primary corpus consists of GATE papers for Computer Science and Information Technology (CS) and for Data Science and
Artificial Intelligence (DA). Table~\ref{tab:sources} lists the sources. The 2025 and 2026 papers and keys were
downloaded from the official websites of the organising institutes (IIT Roorkee and IIT Guwahati). Papers from 2011 to
2024 were obtained from a public mirror; its trustworthiness was established by comparing its 2025 files with the
official ones---the DA paper was byte-identical and the CS paper and key were text-identical. The official GATE 2024 DA
sample paper (IISc Bangalore) supplements the young DA paper~\cite{gate2025,gate2026}. A university end-semester Mathematics-II paper exercises
the non-GATE path.

\begin{table}[h]
\centering
\caption{Corpus sources and provenance}
\label{tab:sources}
\small
\begin{tabularx}{\textwidth}{@{}lXl@{}}
\toprule
\textbf{Years} & \textbf{Papers} & \textbf{Source}\\
\midrule
2025 & CS1, CS2, DA with keys & Official (IIT Roorkee)\\
2026 & CS1, CS2, DA with keys & Official (IIT Guwahati)\\
2011--2024 & CS (1--3 sessions/year), DA 2024; keys from 2017 & Verified mirror\\
2024 & DA sample paper (subject section) & Official (IISc)\\
2025 & Mathematics-II end-semester (university) & Institute paper\\
\bottomrule
\end{tabularx}
\end{table}

Data-quality problems were found and corrected rather than ignored: the mirror's ``2017 CS1'' key was in fact the
Electronics (EC) key and is rejected automatically because it contains no CS rows; the 2013 PDF contained four booklet
codes of the same paper (only one is kept); the 2014, 2015 and 2021 PDFs each bundled two or three different session
papers, which were split into separate papers. The GATE pattern used for slot inference is summarised in
Table~\ref{tab:pattern}.

\begin{table}[h]
\centering
\caption{GATE paper pattern (2014 onwards)}
\label{tab:pattern}
\small
\begin{tabular}{@{}llccl@{}}
\toprule
\textbf{Section} & \textbf{Questions} & \textbf{Marks each} & \textbf{Total} & \textbf{Types}\\
\midrule
General Aptitude & Q1--Q5 & 1 & 5 & MCQ\\
General Aptitude & Q6--Q10 & 2 & 10 & MCQ\\
Subject (CS / DA) & Q11--Q35 & 1 & 25 & MCQ, MSQ, NAT\\
Subject (CS / DA) & Q36--Q65 & 2 & 60 & MCQ, MSQ, NAT\\
\midrule
\multicolumn{3}{@{}l}{\textbf{Total}: 65 questions, 3 hours} & \textbf{100} & \\
\bottomrule
\end{tabular}
\end{table}

\section{Ingestion and Preprocessing}
Figure~\ref{fig:ingest} shows the ingestion decision flow. A document is first read through its PDF text layer. If the
layer is absent (scanned papers) or garbled---detected when more than a quarter of lines are one or two characters long,
the signature of exploded equations---the page images are passed to on-device OCR (Apple Vision); boxes are re-assembled
into lines by their vertical centres so that margin labels such as ``Q.7'' rejoin their text. Multimodal LLM
transcription was evaluated but rejected for GATE papers: the model refuses to reproduce widely published text
(\emph{recitation} filter), which is itself an instructive finding.

\begin{figure}[t]
\centering
\begin{tikzpicture}[node distance=0.55cm]
  \node[blk=3cm] (in) {PDF / images};
  \node[blk=3.2cm, below=of in] (tl) {Read text layer};
  \node[draw, diamond, aspect=2.4, align=center, font=\small, inner sep=1pt, below=of tl] (ok) {usable?};
  \node[blk=2.8cm, right=1.2cm of ok] (ocr) {On-device OCR\\line re-assembly};
  \node[draw, diamond, aspect=2.4, align=center, font=\small, inner sep=1pt, below=of ok] (gate) {GATE layout\\clean?};
  \node[blk=3cm, left=1.1cm of gate] (rx) {Rule parser +\\official key join};
  \node[blk=3cm, right=1.1cm of gate] (llm) {LLM structuring\\(new formats)};
  \node[blk=3.2cm, below=1.3cm of gate] (lab) {Label: topic, Bloom\\$\rightarrow$ store};
  \draw[arr] (in) -- (tl); \draw[arr] (tl) -- (ok);
  \draw[arr] (ok) -- node[above, font=\scriptsize]{no} (ocr);
  \draw[arr] (ocr) |- (gate);
  \draw[arr] (ok) -- node[right, font=\scriptsize]{yes} (gate);
  \draw[arr] (gate) -- node[above, font=\scriptsize]{yes} (rx);
  \draw[arr] (gate) -- node[above, font=\scriptsize]{no} (llm);
  \draw[arr] (rx) |- (lab); \draw[arr] (llm) |- (lab);
\end{tikzpicture}
\caption{Ingestion flow for a new paper}
\label{fig:ingest}
\end{figure}

\subsection{GATE paper parsing and key join}
GATE papers are parsed with layout rules (Algorithm~\ref{alg:parse}). Header, footer and section-banner lines are removed
by a noise pattern. A question begins at a line matching \texttt{Q.\,$n$}; the online-test layout of 2017--2021
(``Question Number : $n$'') and the 2020 variant (``Q.NO.\,$n$'') are normalised to the same form. Options \texttt{(A)}--%
\texttt{(D)} are collected, including options printed inline on one line. Older papers restart numbering after the
aptitude section, so numbers are mapped to a global sequence, and OCR misreadings (``Q.28'' read as ``88'') are repaired
by enforcing monotonic increase. The official key PDF supplies each question's type, section, marks and answer; three
key layouts (2017--2026) are supported, and a key is rejected if it contains no subject rows.

\begin{algorithm}[t]
\caption{GATE paper parsing and answer-key join}
\label{alg:parse}
\begin{algorithmic}[1]
\Require page texts $T$, optional key pages $K$, paper year $y$
\State $L \gets$ normalise each line of $T$; drop lines matching the noise pattern
\State $Q \gets [\,]$
\For{each line $\ell \in L$}
  \If{$\ell$ matches \texttt{Q.}$n$} \State append new question $(n, \text{stem}=[\,], \text{options}=\{\})$ to $Q$
  \ElsIf{$\ell$ matches \texttt{(A)}--\texttt{(D)}} \State start option on the current question
  \Else \State append $\ell$ to the open option or to the stem
  \EndIf
\EndFor
\State renumber $Q$ to a global $1..65$ sequence; repair non-monotonic numbers
\State $\kappa \gets$ \Call{ParseKey}{$K$}; \textbf{if} $\kappa$ has no CS/DA rows \textbf{then} $\kappa \gets \emptyset$
\For{each $q \in Q$}
  \State $(q.\text{type}, q.\text{section}, q.\text{marks}, q.\text{answer}) \gets \kappa[q.n]$ or the default layout for year $y$
\EndFor
\State \Return $Q$
\end{algorithmic}
\end{algorithm}

\subsection{Pattern-agnostic structuring of new papers}
When a paper does not fit the GATE layout---a university paper or a future GATE paper with a changed pattern---its text
is split into chunks and an LLM returns structured questions (number, text in LaTeX, options, marks, section) together
with the list of question formats it observed, so that new formats such as match-the-following or assertion--reason are
captured rather than lost. Offline, a heuristic segmenter recognises question numbers, sub-parts and
``($a\times b$)'' marks headers.

\section{NLP Analysis Layer}
\subsection{Bloom's taxonomy classification}
Each question is assigned a Bloom level by transparent lexical rules, following the finding that directive verbs carry
most of the signal~\cite{mohammed2020}. Cue patterns are checked from the highest level down and the first match wins
(Table~\ref{tab:bloom}); questions with no cue default to \emph{Understand}.

\begin{table}[h]
\centering
\caption{Bloom's taxonomy cue patterns (highest level checked first)}
\label{tab:bloom}
\small
\begin{tabularx}{\textwidth}{@{}lX@{}}
\toprule
\textbf{Level} & \textbf{Example cues}\\
\midrule
Create & design, devise, formulate, construct an algorithm/grammar/automaton, write a program/query\\
Evaluate & justify, evaluate, which is best/most appropriate, assess, recommend\\
Analyze & prove, derive, analyse, compare, time complexity, worst case, which of the following is TRUE/FALSE\\
Apply & calculate, compute, find, determine, solve, number of, value of, output of, probability\\
Understand & explain, describe, discuss, illustrate, interpret\\
Remember & define, state, list, name, what is\\
\bottomrule
\end{tabularx}
\end{table}

\subsection{Topic labelling by distillation}
Each syllabus (official GATE CS and DA syllabi, and the university syllabus) is a list of units, each with topics. A
question must be mapped to one (unit, topic) pair. Labelling proceeds in three stages (Figure~\ref{fig:distill}):
\begin{enumerate}
  \item \textbf{LLM teacher.} Batches of 33 questions are sent to Gemini together with the syllabus; the model returns a
        unit and topic per question, which are snapped to the closest valid syllabus string by sequence similarity.
  \item \textbf{Zero-shot fallback.} Questions the teacher could not label are matched to topic descriptions by cosine
        similarity of MiniLM sentence embeddings~\cite{reimers2019,wang2020minilm}.
  \item \textbf{Student model.} A TF-IDF (word uni- and bi-grams, sublinear tf) + logistic-regression classifier is
        trained on the teacher's labels; the latest year is held out for evaluation. This 7\,MB model labels papers
        uploaded to the deployed server without an LLM call or a deep-learning runtime.
\end{enumerate}
To prevent the student learning from its own predictions, corpus ingestion always uses the teacher, never the student.

\begin{figure}[h]
\centering
\begin{tikzpicture}[node distance=0.5cm and 0.9cm]
  \node[blk=2.3cm] (q) {Corpus\\questions};
  \node[blk=2.6cm, right=of q, fill=blue!6] (t) {LLM teacher\\(+ zero-shot)};
  \node[blk=2.3cm, right=of t] (l) {Labelled\\corpus};
  \node[blk=2.6cm, right=of l] (s) {TF-IDF + LR\\student};
  \node[blk=2.6cm, below=0.8cm of s] (n) {New uploads\\(server)};
  \node[blk=2.3cm, below=0.8cm of l] (e) {Held-out year\\evaluation};
  \draw[arr] (q) -- (t); \draw[arr] (t) -- (l); \draw[arr] (l) -- node[above,font=\scriptsize]{train} (s);
  \draw[arr] (s) -- (n); \draw[arr] (s) -- (e);
\end{tikzpicture}
\caption{Teacher--student labelling scheme}
\label{fig:distill}
\end{figure}

\subsection{Text representation and similarity}
For similarity, each question $q$ (stem plus options; stems alone are often generic, e.g.\ ``Which of the following
statements is FALSE?'') is represented by a TF-IDF vector over character 3--5-grams, which is robust to OCR noise and
inflection:
\begin{equation}
  w_{t,q} = \bigl(1+\log \mathrm{tf}(t,q)\bigr)\cdot \mathrm{idf}(t), \qquad
  \mathrm{idf}(t)=\log\frac{1+|\mathcal{C}|}{1+\mathrm{df}(t)}+1 ,
  \label{eq:tfidf}
\end{equation}
\begin{equation}
  \mathrm{sim}(a,b)=\frac{\mathbf{v}_a\cdot\mathbf{v}_b}{\lVert\mathbf{v}_a\rVert\,\lVert\mathbf{v}_b\rVert}.
  \label{eq:cos}
\end{equation}

\subsection{Historical statistics}
For each unit $u$ the average marks per paper and the share of total marks are computed. The trend is the least-squares
slope $\beta$ of per-paper marks against year index; $\beta>0.5$ is labelled rising, $\beta<-0.5$ falling, otherwise
stable. Topic importance combines marks, frequency and recency in a \emph{hot score}:
\begin{equation}
  h(\tau)=\sum_{q\in\tau}\frac{m_q\,\lambda^{\,Y-y_q}}{n_{y_q}}, \qquad
  H(\tau)=100\cdot\frac{h(\tau)}{\max_{\tau'}h(\tau')},
  \label{eq:hot}
\end{equation}
where $\lambda=0.8$ discounts older papers and division by $n_{y}$ prevents years with several sessions from dominating.
Recurring concepts are pairs of questions from different years with $\mathrm{sim}\ge 0.55$.

\section{Generative Question Paper Pipeline}
Figure~\ref{fig:genpipe} shows the generation pipeline.

\begin{figure}[t]
\centering
\begin{tikzpicture}[node distance=0.45cm]
  \node[blk=4.6cm] (r) {Student request (free text)};
  \node[blk=4.6cm, below=of r] (p) {Parse $\rightarrow$ exam, area, units, topics};
  \node[blk=4.6cm, below=of p] (b) {Blueprint: slots (section, marks, type)};
  \node[blk=4.6cm, below=of b] (pl) {Plan: sample topic per slot by $H(\tau)$};
  \node[blk=4.6cm, below=of pl] (rag) {Retrieve past questions on topic};
  \node[blk=4.6cm, below=of rag, fill=blue!6] (g) {Gemini writes questions + keys};
  \node[blk=4.6cm, below=of g] (v) {Validate format $\cdot$ novelty};
  \node[blk=4.6cm, below=of v] (k) {Blind re-solve of answer keys};
  \node[blk=4.6cm, below=of k] (o) {Paper + validation report};
  \foreach \a/\b in {r/p,p/b,b/pl,pl/rag,rag/g,g/v,v/k,k/o} \draw[arr] (\a) -- (\b);
  \node[blk=3.2cm, right=1.2cm of v] (re) {Regenerate once};
  \node[blk=3.2cm, right=1.2cm of k] (py) {Fallback: real PYQ\\with official key};
  \draw[arr] (v.east) -- node[above,font=\scriptsize]{too similar} (re.west);
  \draw[arr] (k.east) -- node[above,font=\scriptsize]{mismatch} (py.west);
  \draw[arr] (re.north) |- (g.east);
\end{tikzpicture}
\caption{Generative question paper pipeline}
\label{fig:genpipe}
\end{figure}

\subsection{Request understanding}
The request is parsed into a structured \texttt{ExamRequest}---area (\texttt{gate\_cs}, \texttt{gate\_da}, a university
syllabus), exam type (full GATE mock, short practice set, mid-semester or end-semester), focus units, focus topics,
number of questions and difficulty. The LLM receives the full syllabus catalogue and must copy names exactly; outputs
are validated against the syllabus. A regular-expression parser with subject aliases (``OS'', ``DBMS'', ``TOC'') serves
offline. University subjects map onto the GATE syllabus unit of the same name, so a ``mid-sem paper for Operating
Systems'' is grounded in sixteen years of GATE operating-systems questions.

\subsection{Blueprint inference}
A full GATE mock reproduces the slot layout of Table~\ref{tab:pattern}; the split of each marks band between MCQ, MSQ and
NAT is the historical average from the answer keys, rounded with the largest-remainder method so that counts sum to 25
and 30. A mid-semester paper has 5 two-mark questions (all compulsory) and 5 five-mark questions (any four), 30 marks
in 90 minutes, restricted to the first half of each focus unit's topics; an end-semester paper has three sections
totalling 100 marks.

\subsection{History-weighted topic planning}
Each slot is assigned a topic by sampling in proportion to historical importance (Algorithm~\ref{alg:plan}):
\begin{equation}
  P(\tau)\propto H(\tau)+10 ,
  \label{eq:sample}
\end{equation}
where the constant keeps never-asked syllabus topics possible. After a topic is chosen its weight is halved to spread
coverage. Bloom levels are sampled from the corpus distribution, biased towards higher levels for long questions.

\begin{algorithm}[t]
\caption{History-weighted planning of a paper}
\label{alg:plan}
\begin{algorithmic}[1]
\Require slots $S$ from the blueprint, topics in scope $\mathcal{T}$, hot scores $H$, Bloom distribution $B$
\State $w_\tau \gets H(\tau)+10$ for all $\tau\in\mathcal{T}$
\For{each slot $s\in S$}
  \If{$s$ is a General Aptitude slot} \State $s.\tau \gets$ uniform choice of an aptitude topic
  \Else
    \State $s.\tau \gets$ sample $\tau$ with probability $w_\tau/\sum w$; \quad $w_{s.\tau}\gets w_{s.\tau}/2$
  \EndIf
  \State $s.\text{bloom} \gets$ sample from $B$ (higher levels for long questions)
\EndFor
\State \Return $S$
\end{algorithmic}
\end{algorithm}

\subsection{Retrieval-augmented generation}
Slots are grouped by unit into batches of six and generated in parallel. For each slot the prompt contains its type,
marks, unit, topic, Bloom level and difficulty, followed by up to three past questions on the same topic (preferring the
same type and recent years) retrieved from the corpus~\cite{lewis2020}. The system prompt casts the model as a senior
paper-setter, forbids copying references and requires a single verifiable answer with a worked explanation. Output is
constrained to a JSON schema; mathematics is written in LaTeX.

\subsection{Validation and answer-key verification}
Each question is checked for its format (exactly four options for MCQ/MSQ, option letters as the key, a numeric key or
range for NAT). Novelty is measured as the maximum similarity (Eq.~\ref{eq:cos}) to any corpus question; drafts with
$\mathrm{sim}\ge\theta=0.8$ are regenerated with the offending drafts shown as negative examples. Because LLM answer
keys can be wrong~\cite{ji2023}, every objective question is then solved again by an independent call that does not see
the key and may flag the question as ill-posed (Algorithm~\ref{alg:verify}), in the spirit of
self-consistency~\cite{wang2023sc}. A question whose key disagrees is regenerated once and re-verified; if it still fails
it is replaced by a real past question with an official key. Finally the marks of the attemptable questions are summed
and compared with the paper's total.

\begin{algorithm}[t]
\caption{Blind answer-key verification}
\label{alg:verify}
\begin{algorithmic}[1]
\Require generated objective questions $G$ with keys
\State $\hat{a} \gets$ \Call{LLMSolve}{$G$ without keys} \Comment{returns answer and well-posed flag per question}
\State $F \gets \{q\in G : \neg\,\text{wellposed}(q) \lor \neg\,\text{match}(\hat{a}_q, \text{key}_q)\}$
\Comment{NAT: within the key range $\pm1\%$}
\State $G' \gets$ \Call{Regenerate}{$F$}; \quad $F' \gets$ \Call{Verify}{$G'$}
\For{each $q\in F$}
  \If{$q'\in G'\setminus F'$} \State replace $q$ by $q'$ \Else \State replace $q$ by the best-matching real PYQ \EndIf
\EndFor
\State \Return $G$, $|F|$
\end{algorithmic}
\end{algorithm}

\section{Personal Study Assistant}
The assistant is a Gemini chat agent whose tools are PaperMind's own functions (Figure~\ref{fig:agent}):
\texttt{syllabus\_overview}, \texttt{hot\_topics}, \texttt{search\_past\_questions}, \texttt{create\_practice\_set},
\texttt{get\_answer}, \texttt{record\_result} and \texttt{my\_progress}. The model decides when to call a tool, in the
manner of ReAct~\cite{yao2023}; tool schemas are derived from Python type hints and docstrings. Its instructions require
claims about the examination to be grounded in tool output, questions to be asked one at a time, the official key to be
fetched only after the student answers, and every graded attempt to be recorded. Conversation history, including tool
calls, is persisted per student.

For a student, the smoothed accuracy on topic $\tau$ after $c$ correct answers out of $n$ attempts and the resulting
study priority are
\begin{equation}
  a_\tau=\frac{c+1}{n+2}, \qquad \rho_\tau=(1-a_\tau)\,\bigl(H(\tau)+10\bigr),
  \label{eq:priority}
\end{equation}
so a topic is urgent when it is both weak for the student and heavily examined. Laplace smoothing stops a single wrong
answer from marking a topic as completely unknown.

\begin{figure}[t]
\centering
\begin{tikzpicture}[node distance=0.5cm and 1.1cm]
  \node[blk=2.4cm] (s) {Student};
  \node[blk=2.8cm, right=of s, fill=blue!6] (m) {Gemini agent\\(chat + history)};
  \node[blk=3.3cm, right=of m, text width=3.3cm, align=left, font=\scriptsize] (t) {%
    \texttt{syllabus\_overview}\\\texttt{hot\_topics}\\\texttt{search\_past\_questions}\\\texttt{create\_practice\_set}\\
    \texttt{get\_answer}\\\texttt{record\_result}\\\texttt{my\_progress}};
  \node[blk=2.8cm, below=0.8cm of t] (d) {Corpus, generator,\\progress tables};
  \draw[arr, <->] (s) -- node[above,font=\scriptsize]{message} (m);
  \draw[arr, <->] (m) -- node[above,font=\scriptsize]{calls} (t);
  \draw[arr, <->] (t) -- (d);
\end{tikzpicture}
\caption{Tool-using study assistant}
\label{fig:agent}
\end{figure}

\section{Deployment Architecture}
The application is a FastAPI service that serves the REST API and a single-page web interface from one process, packaged
as a Docker image. Only run-time dependencies are installed in the image: the prebuilt corpus database and the 7\,MB
student classifier ship inside it, while heavy ingest-time components (sentence-embedding model, on-device OCR) stay on
the developer machine. Gemini's free tier caps requests per model per day, so two ordered model chains are configured---a
\emph{bulk} chain of lightweight models for structuring and labelling and a \emph{generation} chain led by the strongest
model---and a model whose daily quota is exhausted or which is temporarily overloaded is skipped automatically. The image
runs on a free cloud web-service tier (512\,MB RAM) and redeploys automatically on every push to the repository.

\section{Chapter Summary}
This chapter described the corpus and its provenance, the ingestion pipeline for digital, scanned and unfamiliar papers,
the NLP layer (Bloom rules, teacher--student topic labelling, TF-IDF similarity and historical statistics), the
generation pipeline with blueprint inference, history-weighted planning, retrieval-augmented prompting and blind answer
verification, the tool-using assistant and its priority model, and the deployment architecture.
